\chapter[Petersen]{Polic\'ias y ladrones sobre gr\'aficas de fichas gr\'aficas de inter\'es}
\label{cap:resultados2}

\section{Petersen}
\label{sec:petersen}

Como un resultado agregado en una gr\'afica de inter\'es particular distinta de
un \'arbol, cambiamos nuestro enfoque a la gr\'afica de Petersen, $P$ (ver
\cref{fig:petersen-grafica}).

\begin{figure}[!htbp]
    \centering
    \scalebox{0.75}{
        \begin{tikzpicture}
            \petersen
        \end{tikzpicture}
    }
    \caption{Gr\'afica de Petersen $P$.}
    \label{fig:petersen-grafica}
\end{figure}

Antes de revisar un caso de mayor complejidad, determinaremos el n\'umero
policiaco de la gr\'afica de Petersen. Enunciamos el siguiente lema.

\begin{lema}
\label{lem:dominante-petersen}
    Si $P$ es la gr\'afica de Petersen, entonces $\gamma(P)=3$.
\end{lema}

\begin{proof}
    N\'otese que la gr\'afica de Petersen tiene orden $10$ y todos sus
    v\'ertices tienen grado $3$ y ning\'un lazo. Por tanto, para cualquier
    $S\subseteq V_P$ tal que $S=\l{x,y}$ con $x$ y $y$ distintos,
    \[
        \abs{N(S)} \leq d_P(x) + d_P(y) = 6 \leq 10
    \]
    As\'i, existe un v\'ertice distinto de $x$ y de $y$ que no es adyacente a
    ninguno de los dos, por lo que $\gamma(P)>2$. Por otro lado, se muestra en
    \cref{fig:petersen-dominante} un conjunto dominante de cardinalidad $3$, concluyendo que
    $\gamma(P)=3$.

    \begin{figure}[!htbp]
        \centering
        \scalebox{0.75}{
            \begin{tikzpicture}
                \petersen
                \node (1) [bola2, fill=red]  at ({(360/5)*0 + 90}:5){};
                \node (2) [bola2, fill=red]  at ({(360/5)*2 + 90}:2.5){};
                \node (3) [bola2, fill=red]  at ({(360/5)*3 + 90}:2.5){};
            \end{tikzpicture}
        }
        \caption{Los v\'ertices en rojo forman un conjunto dominante.}
        \label{fig:petersen-dominante}
    \end{figure}

\end{proof}

As\'i, del resultado anterior y d\cref{lem:cotasup-dominante} se obtiene el
siguiente corolario.

\begin{corolario}
\label{cor:cotasup-petersen1}
    Si $P$ es la gr\'afica de Petersen, $c\paren{P}\leq 3$.
\end{corolario}
Luego, tiene lugar el siguiente teorema.

\begin{teorema}
\label{teo:numero-policiaco-petersen}
    Si $P$ denota la gr\'afica de Petersen, $c\paren{P}=3$.
\end{teorema}

\begin{proof}
    En virtud d\cref{cor:cotasup-petersen1}, basta demostrar que $R$ tiene una
    estrategia de escape en un juego de polic\'ias y ladrones sobre $P$ contra
    dos polic\'ias. V\'ease que por lo mencionado en
    \cref{lem:dominante-petersen}, para cualquier par de v\'ertices de $P$ existe
    al menos un v\'ertice a distancia al menos $2$, por lo que $R$ puede escoger
    su posici\'on inicial y sobrevive al primer turno de $C$.
    
    Posterior al turno de $C$, si ning\'un polic\'ia se encuentra a un turno de
    capturar a $r$, $R$ pasa. De lo contrario, sup\'ongase sin p\'erdida de
    generalidad que solo $c_1$ est\'a a un turno de capturarle, de modo que $r$
    tiene a lo m\'as dos posibles movimientos. N\'otese que $N_P(c_1(t))\cap
    N_P(r(t))$ es de cardinalidad uno ya que de lo contrario habr\'ia un ciclo
    de longitud $3$ en $P$, pero la longitud m\'inima de un ciclo en $P$ es $5$.
    As\'i, $c_1$ no restringe ninguno de los otros dos movimientos para $r$.
    Similarmente, $N_P(c_2(t))\cap N_P(r(t))$ no puede ser de cardinalidad $2$
    pues de lo contrario, $P$ tendr\'ia un ciclo de longitud $4$. As\'i, $r$
    tendr\'ia al menos un v\'ertice al cual moverse para sobrevivir al siguiente
    turno de $C$.

    Finalmente, si dos polic\'ias est\'an a un turno de capturar a $r$, por lo
    dicho anteriormente se tiene que ninguno de ellos es adyacente a un vecino
    de $r(t)$. Por tanto, $r$ tiene al menos un movimiento disponible tal que
    sobrevive al siguiente turno de $C$. Esta estrategia se puede iterar siempre
    que sea necesario, por lo que $2<c(P)$ y, consecuentemente, $c(P)=3$.

\end{proof}

Cambiando nuestra atenci\'on a la gr\'afica de $2$ fichas de Petersen,
observemos que d\cref{lem:dominante-petersen} se tiene que la distancia m\'axima
entre dos v\'ertices de $P$ es $2$. Una observaci\'on \'util acerca de $P$ que
se deduce de \cref{fig:petersen-dominante} es que cualesquiera dos de sus
v\'ertices est\'an a lo m\'as a distancia $2$. Adem\'as, dado que $P$ no tiene
ciclos de longitud $4$, sucede que para cualesquiera dos v\'ertices no
adyacentes, $u$ y $v$, son tales que existe una \'unica $uv$-trayectoria de
longitud $2$. As\'i, proponemos el siguiente lema.

% \begin{lema}
% \label{lem:bound}
%       Si $P$ es la gr\'afica de Petersen, $c\paren{\t{2}{P}} > 3$.
% \end{lema}
    
% \begin{proof}
%     Basta probar que $R$ tiene una estrategia de escape contra $3$ polic\'ias.
%     Se demostrar\'a que $R$ es capaz de permanecer en una de dos situaciones:
%     mantiene sus fichas en v\'ertices no adyacentes de $P$ mientras que $C$
%     est\'a a al menos dos turnos de capturarle, o $R$ puede regresar a la
%     primera situaci\'on si es que tiene fichas en v\'ertices adyacentes.

%     N\'otese primero que $C$ puede ocupar, con cada ficha de sus polic\'ias,
%     hasta $6$ v\'ertices. Dado que $P$ tiene orden $10$, $R$ puede escoger su
%     configuraci\'on inicial de entre los $4$ v\'ertices restantes de manera que
%     ninguna ficha de $r$ es capturada por alguna ficha de alg\'un polic\'ia. Por
%     tanto $C$ est\'a a al menos dos turnos de capturarle (v\'ease
%     \cref{fig:initial0}). Adem\'as, dado que $P$ no tiene ciclos de longitud
%     $4$, este tampoco contiene a $K_4$, de manera que al menos dos v\'ertices de
%     los $4$ disponibles son no adyacentes y estos son sobre los que se posiciona
%     $r$.

%     \begin{figure}[!htbp]
%         \begin{tcbitemize}[raster columns=2, raster equal height, enhanced,
%             colback=Black!5!white,colframe=Grey!75!black, attach title to upper,
%             phantom={\setcounter{myrasternum}{\thetcbrasternum}},
%             title={(\alph{myrasternum})\enspace}, halign lower=center, halign
%             upper = left, lower separated=false, middle=0mm, boxsep=0mm,
%             leftupper=2mm, left=2mm, right=2mm, fonttitle=\bfseries,
%             coltitle=black, raster column skip=1mm,raster halign=center, raster
%             row skip=1mm]

%             \tcbitem
%                 \tcblower
%                     \scalebox{0.4}{
%                         \begin{tikzpicture}
%                                 \petersen
%                                 \unopB{5cm}{90}
%                                 \unopG{5cm}{90 + (360/5)*1}
%                                 \unopR{5cm}{90 + (360/5)*2}
%                                 \unopB{5cm}{90 + (360/5)*3}
%                                 \unopG{5cm}{90 + (360/5)*4}
%                                 \unopR{2.5cm}{90}
%                         \end{tikzpicture}
%                     }
            
%             \tcbitem
%                 \tcblower
%                     \scalebox{0.4}{
%                         \begin{tikzpicture}
%                             \petersen
%                             \unopB{5cm}{90}
%                             \unopG{5cm}{90 + (360/5)*1}
%                             \unopR{5cm}{90 + (360/5)*2}
%                             \unopB{5cm}{90 + (360/5)*3}
%                             \unopG{5cm}{90 + (360/5)*4}
%                             \unopR{2.5cm}{90}
%                             \punto{2.5cm}{90 + (360/5)}
%                             \punto{2.5cm}{90 + (360/5)*2}
%                         \end{tikzpicture}
%                     }
%         \end{tcbitemize}
%     \captionsetup{font=scriptsize}
%         \caption{En (a), $C$ ha escogido una posici\'on inicial tal que ocupa la
%         mayor cantidad posible de v\'ertices de $P$. A\'un as\'i, $R$ puede
%         posicionar las fichas de $r$ sobre los v\'ertices restantes como en (b),
%         cumpliendo la condici\'on que busca. N\'otese que todas las fichas
%         tienen la misma forma pues no se utilizar\'a etiquetado alguno en la
%         prueba.}
%     \label{fig:initial0}
%     \end{figure}

%     En los siguientes turnos, $R$ pasa siempre que ning\'un polic\'ia est\'e a
%     un turno de capturarle. Cuando ello pase, obs\'ervese que $R$ tiene $6$
%     posibles movimientos, dos de los cuales le llevan a que sus fichas est\'en
%     en v\'ertices adyacentes debido a que existe una \'unica trayectoria de
%     longitud $2$ entre las posiciones de sus fichas. Sup\'ongase sin p\'erdida
%     de generalidad que $c_1$ est\'a a un turno de capturar a $r$, implicando
%     que, sin p\'erdida de generalidad, $c_{1,1}$ captura a $r_1$ y $c_{1,2}$
%     est\'a a distancia $1$ de $r_2$. Esto implica que una de las fichas de $R$
%     tiene un poisble movimiento menos, a saber, el que lleva $r_2$ al v\'ertice
%     en el que se encuentra $c_{1,2}$. M\'as a\'un, obs\'ervese que cualquier
%     polic\'ia a al menos $3$ turnos de capturar a $R$ no representa un problema
%     pues el movimiento de $R$ le llevar\'a a estar a al menos $2$ turnos de
%     capturarle, por lo que $R$ sobrevive al siguiente turno de dicho polic\'ia.
%     As\'i, solo se considerar\'an los polic\'ias a distancia $1$ o $2$ de
%     capturar a $r$. Si $c_1$ es el \'unico polic\'ia a menos de $3$ turnos de
%     capturar a $r$, entonces $R$ puede mover $r_1$ a un v\'ertice no adyacente a
%     la posici\'on de $r_2$, consiguiendo su objetivo (ver
%     \cref{fig:best_cover1}).

%     \begin{figure}[!htbp]
%         \begin{tcbitemize}[raster columns=2, raster equal height, enhanced,
%             colback=Black!5!white,colframe=Grey!75!black, attach title to upper,
%             phantom={\setcounter{myrasternum}{\thetcbrasternum}},
%             title={(\alph{myrasternum})\enspace}, halign lower=center, halign
%             upper = left, lower separated=false, middle=0mm, boxsep=0mm,
%             leftupper=2mm, left=2mm, right=2mm, fonttitle=\bfseries,
%             coltitle=black, raster column skip=1mm,raster halign=center, raster
%             row skip=1mm]
    
%             \tcbitem
%                 \tcblower
%                     \scalebox{0.4}{
%                         \begin{tikzpicture}
%                             \localpetersen
%                             \begin{scope}[xshift=2.5cm]
%                                 \unopB{2.5cm}{(360/3)*2+180}
%                             \end{scope}
%                             \unopBp{2.5cm}{180}
%                             \punto{2.5cm}{0}
%                         \end{tikzpicture}
%                     }
            
%             \tcbitem
%                 \tcblower
%                     \scalebox{0.4}{
%                         \begin{tikzpicture}
%                             \localpetersen
%                             \begin{scope}[xshift=2.5cm]
%                                 \unopB{2.5cm}{(360/3)*2 + 180}
%                             \end{scope}
%                             \unopB{2.5cm}{180}
%                             \begin{scope}[xshift=-2.5cm]
%                                 \punto{2.5cm}{(360/3)*1}
%                             \end{scope}
%                             \punto{2.5cm}{0}
%                         \end{tikzpicture}
%                     }
%         \end{tcbitemize}
%         \captionsetup{font=scriptsize}
%         \caption{En (a) se muestran los v\'ertice a los que $R$ puede mover una
%         de las fichas de $r$. Adem\'as, $c_1$ es el \'unico a menos de $3$
%         turnos de capturar a $r$. Incluso si $c_2$ y $c_3$ tuvieran algunas de
%         sus fichas en las vecindades de $r_1(t)$ y $r_2(t)$, el movimiento de
%         $r_1$ hacia alguno de esos v\'ertices no lo acercan lo suficiente a
%         $c_2$ y $c_3$, de manera que $c_2$ y $c_3$ no capturan a $r$ en su
%         siguiente turno. As\'i, $R$ tiene $3$ movimientos que eviten que alg\'un
%         polic\'ia de $C$ le capture al siguiente turno. Uno de dichos
%         movimientos se muestra en (b).}
%         \label{fig:best_cover1}
%     \end{figure}
    
%     Similarmente y sin p\'erdida de generalidad, sup\'ongase que $c_1$ y $c_2$
%     son los \'unicos polic\'ias a menos de $3$ turnos de capturar a $r$. Se
%     generan varios casos dependiendo de la distancia entre las fichas de $c_2$ y
%     las fichas de $r$. 
    
%     \begin{enumerate}
%         \item Si $c_2$ tambi\'en est\'a a un turno de capturar a $r$ entonces
%         este prohibe a lo m\'as un v\'ertice adicional de los $3$ disponibles
%         mencionados en \cref{fig:best_cover1} (v\'ease (a) y (b) en
%         \cref{fig:best_cover2}).
        
%         \item Si $c_2$ est\'a a dos turnos de capturar a $r$, se generan dos
%         casos adicionales.
%         \begin{enumerate}
%             \item Si $c_2$ captura una ficha de $r$ ($r_1$ capturada por
%             $c_{2,1}$ sin p\'erdida de generalidad) y $r_2$ est\'a a distancia
%             $2$ de $c_{2,2}$, al igual que en (a) de \cref{fig:best_cover2},
%             $c_2$ puede prohibir a lo m\'as un movimiento m\'as de $R$ ya que,
%             de lo contrario, habr\'ia un v\'ertice adyacente a $2$ vecinos de
%             $r_2(t)$, lo cual es imposible pues $P$ no tiene ciclos de longitud
%             $4$ (v\'ease (c) y (d) en \cref{fig:best_cover2}).

%             \item Si ambos $c_{2,1}$ y $c_{2,2}$ est\'an, sin p\'erdida de
%             generalidad, a distancia $1$ de $r_1$ y $r_2$ respectivamente,
%             entonces $c_2$ prohibe a lo m\'as dos movimientos adicionales de
%             $R$, uno por cada ficha de $r$. As\'i, de los $3$ posibles
%             movimientos en \cref{fig:best_cover1}, al menos uno de ellos sigue
%             disponible para $R$ (ver (e) y (f) en \cref{fig:best_cover2}).
%         \end{enumerate}
%         N\'otese que en efecto estos son los \'unicos dos casos pues la \'unica
%         otra manera en que $c_2$ est\'e a dos turnos de capturar a $r$ es que
%         dicho polic\'ia capture a una ficha de $r$ con, suponidendo sin
%         p\'erdida de generalidad, $c_{2,1}$, y $c_{2,2}$ est\'e obstruida por
%         $c_{2,1}$. No obstante, dado que la distancia entre cualesquiera dos
%         v\'ertices de $P$ es a lo m\'as $2$, esto implicar\'ia que ambas fichas
%         de $R$ est\'an en v\'ertices adyacentes, lo cual no es el caso.
%     \end{enumerate}
    
%     \begin{figure}[!htbp]
%         \begin{tcbitemize}[raster columns=2, raster equal height, enhanced,
%             colback=Black!5!white,colframe=Grey!75!black, attach title to upper,
%             phantom={\setcounter{myrasternum}{\thetcbrasternum}},
%             title={(\alph{myrasternum})\enspace}, halign lower=center, halign
%             upper = left, lower separated=false, middle=0mm, boxsep=0mm,
%             leftupper=2mm, left=2mm, right=2mm, fonttitle=\bfseries,
%             coltitle=black, raster column skip=1mm,raster halign=center, raster
%             row skip=1mm]
    
%             \tcbitem
%                 \tcblower
%                     \scalebox{0.4}{
%                         \begin{tikzpicture}
%                             \localpetersen
%                             \begin{scope}[xshift=2.5cm]
%                                 \unopB{2.5cm}{(360/3)*2 + 180}
%                             \end{scope}
%                             \unopBp{2.5cm}{180}
%                             \unopGp{2.5cm}{0}
%                             \begin{scope}[xshift=-2.5cm]
%                                 \unopG{2.5cm}{(360/3)*1}
%                             \end{scope}
%                         \end{tikzpicture}
%                     }
            
%             \tcbitem
%                 \tcblower
%                     \scalebox{0.4}{
%                         \begin{tikzpicture}
%                             \localpetersen
%                             \begin{scope}[xshift=2.5cm]
%                                 \unopB{2.5cm}{(360/3)*2 + 180}
%                             \end{scope}
%                             \unopB{2.5cm}{180}
%                             \unopGp{2.5cm}{0}
%                             \begin{scope}[xshift=-2.5cm]
%                                 \unopG{2.5cm}{(360/3)*1}
%                                 \punto{2.5cm}{(360/3)*2}
%                             \end{scope}
%                         \end{tikzpicture}
%                     }
            
%             \tcbitem
%                 \tcblower
%                     \scalebox{0.4}{
%                         \begin{tikzpicture}
%                             \localpetersen
%                             \begin{scope}[xshift=2.5cm]
%                                 \unopB{2.5cm}{(360/3)*2 + 180}
%                                 \unopG{2.5cm}{(360/3)*1 + 180}
%                             \end{scope}
%                             \unopBp{2.5cm}{180}
%                             \unopGp{2.5cm}{0}
%                         \end{tikzpicture}
%                     }
            
%             \tcbitem
%                 \tcblower
%                     \scalebox{0.4}{
%                         \begin{tikzpicture}
%                             \localpetersen
%                             \begin{scope}[xshift=2.5cm]
%                                 \unopB{2.5cm}{(360/3)*2 + 180}
%                                 \unopG{2.5cm}{(360/3)*1 + 180}
%                             \end{scope}
%                             \unopB{2.5cm}{180}
%                             \unopGp{2.5cm}{0}
%                             \begin{scope}[xshift=-2.5cm]
%                                 \punto{2.5cm}{(360/3)}
%                             \end{scope}
%                         \end{tikzpicture}
%                     }
        
%             \tcbitem
%                 \tcblower
%                     \scalebox{0.4}{
%                         \begin{tikzpicture}
%                             \localpetersen
%                             \begin{scope}[xshift=2.5cm]
%                                 \unopB{2.5cm}{(360/3)*2 + 180}
%                                 \unopG{2.5cm}{(360/3)*1 + 180}
%                             \end{scope}
%                             \begin{scope}[xshift=-2.5cm]
%                                 \unopG{2.5cm}{(360/3)*1}
%                             \end{scope}
%                             \unopBp{2.5cm}{180}
%                             \punto{2.5cm}{0}
%                         \end{tikzpicture}
%                     }
            
%             \tcbitem
%                 \tcblower
%                     \scalebox{0.4}{
%                         \begin{tikzpicture}
%                             \localpetersen
%                             \begin{scope}[xshift=2.5cm]
%                                 \unopB{2.5cm}{(360/3)*2 + 180}
%                                 \unopG{2.5cm}{(360/3)*1 + 180}
%                             \end{scope}
%                             \begin{scope}[xshift=-2.5cm]
%                                 \unopG{2.5cm}{(360/3)*1}
%                                 \punto{2.5cm}{(360/3)*2}
%                             \end{scope}
%                             \unopB{2.5cm}{180}
%                             \punto{2.5cm}{0}
%                         \end{tikzpicture}
%                     }
%         \end{tcbitemize}
%         \captionsetup{font=scriptsize}
%         \caption{En (a), a pesar de que $c_1$ y $c_2$ est\'an cerca de capturar
%         a $r$, este puede moverse como en (b) para escapar de ambos y mantenerse
%         en una posici\'on deseada, notando que $c_3$ se ignora pues no est\'a lo
%         suficientemente cerca de capturar a $r$. Para (c), $R$ a\'un tiene dos
%         movimientos posibles (uno mostrado en (d)) que no pone a sus fichas en
%         v\'ertices adyacentes, mientras que para (e), tiene exactamente un
%         movimiento que lo deja en una posici\'on deseada.}
%         \label{fig:best_cover2}
%     \end{figure}

%     As\'i, resta el caso en el que todos los polic\'ias est\'an a lo m\'as a dos
%     turnos de capturar a $r$. El razonamiento previo implica que si los tres
%     polic\'ias est\'an a un turno de capturar a $r$, entonces este tiene a lo
%     m\'as $3$ de $4$ v\'ertices prohibidos que lo mantienen en una posici\'on
%     deseada. De forma similar, si todos salvo un polic\'ia est\'an a un turno de
%     capturarle, sup\'ongase sin p\'erdida de generalidad que $c_3$ est\'a a dos
%     turnos de capturarle. As\'i $c_1$ y $c_2$ prohiben, entre los dos, a lo
%     m\'as dos de los $4$ posibles movimientos de $R$, mientras que $c_3$ prohibe
%     a lo m\'as uno adicional por lo dicho en el caso 2a), de modo que $R$ a\'un
%     tiene un movimiento que lo deja en una posici\'on deseada. Por tanto, solo
%     resta atender dos casos:
%     \begin{enumerate}
%         \item Dos polic\'ias est\'an a un turno de capturar a $r$ (capturando
%         fichas distintas de $r$) mientras que el polic\'ia restante no captura
%         ninguna ficha de $r$.
        
%         \item Un solo polic\'ia captura una ficha de $R$ mientras que los otros
%         dos cubren las vecindades $r_1$ y $r_2$ (cada polic\'ia tiene una ficha
%         en un vecino de $r_1(t)$ y una ficha en un vecino de $r_2(t)$
%     \end{enumerate}
%     Se resolver\'a \'unicamente el segundo caso pues propicia movimientos
%     forzados para $R$, mientras que el primero es una versi\'on m\'as simple del
%     segundo. As\'i pues, $R$ est\'a obligado a mover una ficha de $r$ hacia el
%     \'unico v\'ertice en $N_P(r_1(t))\cap N_P(r_2(t))$, sobreviviendo al
%     siguiente turno de $C$ pues ning\'un polic\'ia est\'a a un turno de
%     capturarle (v\'ease (b) en \cref{fig:possible}). Sup\'ongase sin p\'erdida
%     de generalidad que mueve $r_1$.


%     \begin{figure}[!htbp]
%     \begin{tcbitemize}[raster columns=2, raster equal height, enhanced,
%         colback=Black!5!white,colframe=Grey!75!black, attach title to upper,
%         phantom={\setcounter{myrasternum}{\thetcbrasternum}},
%         title={(\alph{myrasternum})\enspace}, halign lower=center, halign
%         upper = left, lower separated=false, middle=0mm, boxsep=0mm,
%         leftupper=2mm, left=2mm, right=2mm, fonttitle=\bfseries,
%         coltitle=black, raster column skip=1mm,raster halign=center, raster
%         row skip=1mm]

%         \tcbitem
%             \tcblower
%                 \scalebox{0.4}{
%                     \begin{tikzpicture}
%                     \petersen
%                     \unopB{2.5cm}{90 + (360/5)*4}
%                     \unopG{5cm}{90 + (360/5)*1}
%                     \unopG{5cm}{90 + (360/5)*2}
%                     \unopR{2.5cm}{90 + (360/5)*3}
%                     \unopR{2.5cm}{90}
%                     \unopBp{2.5cm}{90 + (360/5)}
%                     \punto{2.5cm}{90 + (360/5)*2}
%                     \end{tikzpicture}
%                 }
        
%         \tcbitem
%             \tcblower
%                 \scalebox{0.4}{
%                     \begin{tikzpicture}
%                     \petersen
%                     \unopBp{2.5cm}{90 + (360/5)*4}
%                     \unopG{5cm}{90 + (360/5)*1}
%                     \unopG{5cm}{90 + (360/5)*2}
%                     \unopR{2.5cm}{90 + (360/5)*3}
%                     \unopR{2.5cm}{90}
%                     \unopB{2.5cm}{90 + (360/5)}
%                     \punto{2.5cm}{90 + (360/5)*2}
%                     \end{tikzpicture}
%                 }

%         \tcbitem
%             \tcblower
%                 \scalebox{0.4}{
%                     \begin{tikzpicture}
%                     \petersen
%                     \unopB{5cm}{90 + (360/5)*4}
%                     \unopG{5cm}{90 + (360/5)*2}
%                     \unopR{2.5cm}{90 + (360/5)*3}
%                     \unopBunopG{2.5cm}{90 + (360/5)}
%                     \unopRp{2.5cm}{90 + (360/5)*2}
%                     \punto{2.5cm}{90 + (360/5)*4}
%                     \end{tikzpicture}
%                 }
        
%         \tcbitem
%             \tcblower
%                 \scalebox{0.4}{
%                     \begin{tikzpicture}
%                     \petersen
%                     \unopBp{5cm}{90 + (360/5)*4}
%                     \unopG{5cm}{90 + (360/5)*2}
%                     \unopR{2.5cm}{90 + (360/5)*3}
%                     \unopBunopG{2.5cm}{90 + (360/5)}
%                     \unopRp{2.5cm}{90 + (360/5)*2}
%                     \end{tikzpicture}
%                 }
%     \end{tcbitemize}
%     \captionsetup{font=scriptsize}
%     \caption{En (a) se muestra la forma en la que $C$ prohibe la mayor cantidad
%     de movimientos posibles para $R$, de manera que este est\'a obligado a poner
%     las fichas de $r$ en v\'ertices adyacentes. N\'otese que el movimiento de
%     $R$ mostrado en (b) asegura que $r_1$ tiene un v\'ertice adicional al cual
%     moverse (aquel distinto de $r_1(t-1)$ y $r_2(t)$). En (c), los polic\'ias se
%     mueven tratando de acercarse a las fichas de $r$ mientras que estos
%     contin\'uan prohibiendo movimientos para $R$, mas $R$ a\'un puede moverse
%     como en (d) sin acercarse a las fichas de los polic\'ias.}
%     \label{fig:possible}
%     \end{figure}

%     Se afirma que, sin importar que haga $C$, este no logra evitar que $R$ mueva
%     $r_1$ hacia el v\'ertice $u$ distinto de $r_1(t-1)$ y $r_2(t)$. Dicho
%     v\'ertice est\'a a distancia $2$ de $r_2(t)$, por lo que probar el hecho
%     anterior implicar\'ia que $R$ logr\'o mantener a sus fichas en v\'ertices no
%     adyacentes y ning\'un polic\'ia se encuentra a un turno de capturarle. As\'i
%     pues, n\'otese que en $u$ no puede haber fichas de alg\'un polic\'ia m\'as
%     que de $c_1$, lo cual no representa una amenaza para $R$ pues la otra ficha
%     de $c_1$ est\'a a $2$ turnos de capturar a $r_2$. M\'as a\'un, como $c_2$ y
%     $c_3$ est\'an a $3$ turnos de capturar a $r$ despu\'es del movimiento de $R$
%     (como se muestra en (b) de \cref{fig:possible}), entonces en el turno de $C$
%     estos se pueden acercar a estar a $2$ turnos de capturar a $R$. No obstante,
%     cuando $R$ mueve $r_1$ hacia $u$, $r$ preserva su distancia hacia $c_2$ en
%     $F_2(P)$ o la incremente en una unidad (v\'ease (c) y (d) en
%     \cref{fig:possible}). As\'i, $R$ regresa a la misma posici\'on inicial en la
%     que ning\'un polic\'ia est\'a a un turno de capturarle y sus fichas est\'an
%     en v\'ertices no adyacentes, asegurando que tiene una estrategia para
%     impedir que el juego termine. Se anexa en \cref{fig:possible2} el otro caso
%     con un argumento similar para asegurar el escape de $R$. Luego,
%     $c\paren{\t{2}{P}}>3$.

%     \begin{figure}[!htbp]
%         \begin{tcbitemize}[raster columns=2, raster equal height, enhanced,
%             colback=Black!5!white,colframe=Grey!75!black, attach title to upper,
%             phantom={\setcounter{myrasternum}{\thetcbrasternum}},
%             title={(\alph{myrasternum})\enspace}, halign lower=center, halign
%             upper = left, lower separated=false, middle=0mm, boxsep=0mm,
%             leftupper=2mm, left=2mm, right=2mm, fonttitle=\bfseries,
%             coltitle=black, raster column skip=1mm,raster halign=center, raster
%             row skip=1mm]
    
%             \tcbitem
%                 \tcblower
%                     \scalebox{0.4}{
%                         \begin{tikzpicture}
%                             \petersen
%                             \unopG{5cm}{90 + (360/5)*1}
%                             \unopB{5cm}{90 + (360/5)*2}
%                             \unopR{2.5cm}{90 + (360/5)*3}
%                             \unopR{2.5cm}{90}
%                             \unopBp{2.5cm}{90 + (360/5)}
%                         \unopGp{2.5cm}{90 + (360/5)*2}
%                         \end{tikzpicture}
%                     }
            
%             \tcbitem
%                 \tcblower
%                     \scalebox{0.4}{
%                         \begin{tikzpicture}
%                             \petersen
%                             \unopG{5cm}{90 + (360/5)*1}
%                             \unopB{5cm}{90 + (360/5)*2}
%                             \unopR{2.5cm}{90 + (360/5)*3}
%                             \unopR{2.5cm}{90}
%                             \unopBp{2.5cm}{90 + (360/5)}
%                             \unopG{2.5cm}{90 + (360/5)*2}
%                             \punto{2.5cm}{90 + (360/5)*4}
%                         \end{tikzpicture}
%                     }
    
%             \tcbitem
%                 \tcblower
%                     \scalebox{0.4}{
%                         \begin{tikzpicture}
%                             \petersen
%                             \unopG{5cm}{90 + (360/5)*1}
%                             \unopB{5cm}{90 + (360/5)*3}
%                             \unopR{2.5cm}{90 + (360/5)*3}
%                             \unopBp{2.5cm}{90 + (360/5)}
%                             \unopGunopR{2.5cm}{90 + (360/5)*2}
%                             \punto{2.5cm}{90 + (360/5)*4}
%                         \end{tikzpicture}
%                     }
            
%             \tcbitem
%                 \tcblower
%                     \scalebox{0.4}{
%                         \begin{tikzpicture}
%                             \petersen
%                             \unopG{5cm}{90 + (360/5)*1}
%                             \unopB{5cm}{90 + (360/5)*3}
%                             \unopR{2.5cm}{90 + (360/5)*3}
%                             \unopBp{2.5cm}{90 + (360/5)}
%                             \unopGunopR{2.5cm}{90 + (360/5)*2}
%                             \punto{2.5cm}{90 + (360/5)*4}
%                         \end{tikzpicture}
%                     }
%         \end{tcbitemize}
%         \captionsetup{font=scriptsize}
%         \caption{En (a), $R$ est\'a obligado a mover una de sus dos fichas hacia
%         el \'unico v\'ertice en $N_P(r_1(t))\cap N_P(r_2(t))$, sup\'ongase sin
%         p\'erdida de generalidad que mueve $r_2$ como en (b). En (c), $c_1$ se
%         ve obligado a hacer dicho movimiento para que $r_1$ no vaya hacia el
%         v\'ertice $u$ como en (d) de \cref{fig:possible}. Si $c_3$ no mueve,
%         $r_1$ se mueve al v\'ertice adyacente que tiene una ficha de $c_3$,
%         notando que $R$ tendr\'ia sus fichas en v\'ertices no adyacentes y
%         ning\'un polic\'ia le capturar\'ia al siguiente turno sin importar el
%         movimiento de $c_2$. Todos los posibles movimientos de $c_3$, salvo por
%         el mostrado en (c), llevan al mismo resultado. Finalmente, n\'otese que
%         en (c) ning\'un polic\'ia est\'a a un turno de capturar a $r$.}
%         \label{fig:possible2}
%         \end{figure}

% \end{proof}
    
\begin{teorema}
    Si $P$ denota la gr\'afica de Petersen, entonces $c\paren{\t{2}{P}}\leq4$.
\end{teorema}

\begin{proof}
    La estrategia ganadora de $C$ con $4$ polic\'ias se divide en tres etapas:
    \begin{enumerate}
        \item La captura de $r_1$ con $c_{1,1}$ y $c_{3,1}$ y la captura de
        $r_2$ con $c_{2,2}$ y $c_{4,2}$.

        \item El posicionamiento de $r_1$ y $r_2$ en v\'ertices adyacentes.
    
        \item La captura de $R$.
    \end{enumerate}

    Para la primera etapa, como $\gamma(P)=3$, existe $S\subseteq V_P$ de
    cardinalidad $3$ tal que este es un conjunto dominante de $P$. Luego, $C$
    puede escoger su configuraci\'on inicial de manera que, para cada dos
    v\'ertices $x$ y $y$ de $S$, existe $i\in\seg{1}{4}$ tal que $c_i$ tiene sus
    fichas en $x$ y $y$ (v\'ease (a) en \cref{fig:initial}). Se afirma que, en
    su primer turno despu\'es de escoger su configuraci\'on inicial, $c_{1,1}$ y
    $c_{2,2}$ pueden capturar a $r_1$ y $r_2$ respectivamente. Esto se puede ver
    en \cref{fig:initial} pues en cada v\'ertice hay al menos dos colores, de
    modo que si $c_{1,1}$ y $c_{2,2}$ no estuvieran a distancia $1$ de sus
    objetivos, estos solo intercambiar\'ian sus colores (y sus formas en caso de
    ser necesario). Por otro lado, $c_{3,1}$ y $c_{4,2}$ se mueven para capturar
    $r_1$ y $r_2$ respectivamente. As\'i, si $c_3$ tuviera alguna de sus fichas
    a distancia $1$ de $r_1$, este reetiqueta en caso de ser necesario para que
    $c_{3,1}$ capture a $r_1$ en el mismo turno (y lo mismo sucede para
    $c_{4,2}$). No obstante, para cada $i\in\seg{1}{4}$, hay un v\'ertice de $S$
    en el que no hay ninguna ficha de $c_i$, de modo que $c_{3,1}$ y $c_{4,2}$
    puede que no capturen a sus objetivos en el primer turno de $C$. A pesar de
    ello, dado que la distancia entre cada par de v\'ertices de $P$ es a lo
    m\'as $2$, en el peor de los casos ambos $c_{3,1}$ y $c_{4,2}$ se encuentran
    a distancia $1$ de sus respectivos objetivos (v\'ease (b) en
    \cref{fig:initial}).

    \begin{figure}[!htbp]
        \begin{tcbitemize}[raster columns=2, raster equal height, enhanced,
            colback=Black!5!white,colframe=Grey!75!black, attach title to upper,
            phantom={\setcounter{myrasternum}{\thetcbrasternum}},
            title={(\alph{myrasternum})\enspace}, halign lower=center, halign upper
            = left, lower separated=false, middle=0mm, boxsep=0mm, leftupper=2mm, left=2mm,
            right=2mm, fonttitle=\bfseries, coltitle=black, raster column skip=1mm,raster
            halign=center, raster row skip=1mm]

            \tcbitem
                \tcblower
                    \scalebox{0.4}{
                        \begin{tikzpicture}
                                \petersen
                                \unopBunopYunotR{5cm}{90 + (360/5)*2}
                                \unotBunotYunopG{2.5cm}{90 + (360/5)*4}
                                \unotGunopR{2.5cm}{90}
                        \end{tikzpicture}
                    }
            
            \tcbitem
                \tcblower
                    \scalebox{0.4}{
                        \begin{tikzpicture}
                            \petersen
                            \unopBp{5cm}{90 + (360/5)*3}
                            \unopR{2.5cm}{90 + (360/5)*3}
                            \unotGt{5cm}{90}
                            \unotY{5cm}{90 + (360/5)*4}
                            \unotBunopG{2.5cm}{90 + (360/5)*4}
                            \unotRunopY{5cm}{90 + (360/5)*2}
                        \end{tikzpicture}
                    }
            \end{tcbitemize}
        \captionsetup{font=scriptsize }
        \caption{Una posible configuraci\'on inicial se muestra en (a), donde
        todas las fichas yacen en los tres v\'ertices de un conjunto dominante.
        Se fijan colores para diferenciar distintos polic\'ias, de modo que el
        azul corresponde a $c_1$, verde a $c_2$, rojo a $c_3$ y amarillo a
        $c_4$. De forma similar, las figuras corresponden al n\'umero de ficha
        del polic\'ia correspondiente, implicando que el punto azul representa
        la ficha $c_{1,1}$, mientras que el tri\'angulo verde representa la
        ficha $c_{2,2}$. En (b) se muestra una situaci\'on posterior al primer
        turno de $C$.}
        \label{fig:initial}
    \end{figure}

    Si $R$ pasa, ambos $c_{3,1}$ y $c_{4,2}$ capturan sus respectivos objetivos
    al siguiente turno. De otro modo, sup\'ongase sin p\'erdida de generalidad
    que $R$ mueve $r_1$. Esto lleva a que $c_{4,2}$ captura $r_2$, mientras que
    $c_{1,1}$ recaptura $r_1$ y $c_{3,1}$ recorre la trayectoria m\'as corta
    hacia $r_1$, la cual es \'unica pues $P$ no tiene ciclos de longitud $4$.
    As\'i, $c_{3,1}$ vuelve a estar a distancia $1$ de $r_1$. Si $R$ pretende
    evitar que ambos $c_{4,2}$ y $c_{3,1}$ capturen sus objetivos, este tiene
    que mover $r_1$ por el resto de sus turnos. M\'as a\'un, tiene que mover
    $r_1$ hacia v\'ertices fuera de la vecindad cerrada de $c_{3,2}(t)$, pues al
    moverse a alguno de estos, $c_3$ reetiqueta sus fichas de manera que
    $c_{3,1}$ captura a $r_1$. N\'otese que $R$ tendr\'ia que mover $r_1$ sobre
    un $C_6$ inducido en $G$ (sin capaz de moverse mas que en una direcci\'on
    pues $c_{3,1}$ le bloquea un v\'ertice sobre el $C_6$, v\'ease (a) en
    \cref{fig:cycle}). As\'i, $c_2$ o $c_4$ podr\'ian posicionar su primer ficha
    sobre el ciclo de manera que, eventualmente, $r_1$ no puede avanzar sobre el
    $C_6$ pues ambos $c_{2,2}$ y $c_{4,2}$ ya captuan a $r_2$. Por tanto, $R$
    solo puede pasar o mover $r_1$ hacia atr\'as sobre $C_6$ (de manera que
    $c_{3,1}$ captura a $r_1$), o mover $r_1$ hacia la vecindad de $c_{3,2}(t)$.
    En cualquiera de los casos, $c_{3,1}$ termina por capturar a $r_1$, de
    manera que ambos $c_{3,1}$ y $c_{4,2}$ capturan a sus respectivos objetivos
    (v\'ease (d) en \cref{fig:cycle}).

    \begin{figure}[!htbp]
        \begin{tcbitemize}[raster columns=2, raster equal height, enhanced,
            colback=Black!5!white,colframe=Grey!75!black, attach title to upper,
            phantom={\setcounter{myrasternum}{\thetcbrasternum}},
            title={(\alph{myrasternum})\enspace}, halign lower=center, halign
            upper = left, lower separated=false, middle=0mm, boxsep=0mm,
            leftupper=2mm, left=2mm, right=2mm, fonttitle=\bfseries,
            coltitle=black, raster column skip=1mm,raster halign=center, raster
            row skip=1mm]

            \tcbitem
                \tcblower
                    \scalebox{0.4}{
                        \begin{tikzpicture}
                                \petersendos
                                \unotGunotYt{5cm}{90 + (360/5)}
                                \unopBp{2.5cm}{90 + (360/5)*4}
                                \unopR{5cm}{90 + (360/5)*4}
                                \unotRunopY{5cm}{90 + (360/5)*2}
                                \unopG{5cm}{90}
                                \unotB{2.5cm}{90 + (360/5)}

                        \end{tikzpicture}
                    }
            
            \tcbitem
                \tcblower
                    \scalebox{0.4}{
                        \begin{tikzpicture}
                            \petersendos
                            \unotGunotYt{5cm}{90 + (360/5)}
                            \unopBp{2.5cm}{90 + (360/5)}
                            \unotBunopR{2.5cm}{90 + (360/5)*4}
                            \unotRunopY{5cm}{90 + (360/5)*2}
                            \unopG{5cm}{90}
                            
                        \end{tikzpicture}
                    }
            \tcbitem
                \tcblower
                    \scalebox{0.4}{
                        \begin{tikzpicture}
                            \petersendos
                            \unotGunotYt{5cm}{90 + (360/5)}
                            \unopBp{2.5cm}{90 + (360/5)*3}
                            \unotB{2.5cm}{90 + (360/5)*4}
                            \unopR{2.5cm}{90 + (360/5)}
                            \unotRunopY{5cm}{90 + (360/5)*2}
                            \unopG{5cm}{90}
                            
                        \end{tikzpicture}
                    }
            
                \tcbitem
                    \tcblower
                        \scalebox{0.4}{
                            \begin{tikzpicture}
                                \petersendos
                                \unotGunotYt{5cm}{90 + (360/5)}
                                \unopBp{2.5cm}{90}
                                \unotB{2.5cm}{90 + (360/5)*4}
                                \unopR{2.5cm}{90 + (360/5)*3}
                                \unotRunopY{5cm}{90 + (360/5)*2}
                                \unopG{5cm}{90}
                                
                            \end{tikzpicture}
                        }
        \end{tcbitemize}
        \captionsetup{font=scriptsize}
        \caption{Cada escena representa un turno del juego donde $R$ debe mover.
        Obs\'ervese que $r_1$ evita la vecindad de $c_{3,2}(t)$ y escapa de
        $c_{3,1}$ movi\'endose sobre el ciclo morado. Despu\'es de $3$ turnos,
        si $r_1$ se quiere mover, este termina capturado por $c_{3,1}$ o se
        posiciona a turno $1$ de que $c_{3,2}$ le capture. Siempre que se tenga
        una obstrucci\'on, las fichas simplemente se reetiquetan como en (b).}
        \label{fig:cycle}
    \end{figure}

    Para la segunda etapa y en los siguientes turnos, $c_{1,1}$ y $c_{3,1}$
    recapturan a $r_1$ si $R$ lo mueve. De la misma manera, $c_{2,2}$ y
    $c_{4,2}$ recapturan cada que $R$ mueva $r_2$. Sea $N(r_1(t))=\l{u_1,v,u_2}$
    y $N(r_2(t))=\l{w_1,v,w_2}$, notando que la intersecci\'on de ambos
    conjuntos es $\l{v}$ pues entre cualesquiera dos v\'ertices no adyacentes
    existe una \'unica trayctoria de longitud $2$. Sin p\'erdida de generalidad,
    sup\'ongase que la \'ulitma ficha movida por $R$ es $r_1$. Luego, $c_1$ y
    $c_3$ usan su movimiento para recapturar tal como se dijo al inicio del
    p\'arrafo. En tanto, sin p\'erdida de generalidad, $c_{2,1}$ y $c_{4,1}$ se
    mueven hacia un v\'ertice adyacente a $u_1$ y $u_2$ respectivamente.
    N\'otese que esto es posible pues, desde cualquier v\'ertice en el que
    $c_{2,1}$ y $c_{4,1}$ pudieran estar, $C$ puede mover uno hacia la vecindad
    de $u_1$ y el otro a la vecindad de $u_2$ (v\'ease \cref{fig:guarding0}).

    \begin{figure}[!htbp]
        \begin{tcbitemize}[raster columns=2, raster equal height, enhanced,
            colback=Black!5!white,colframe=Grey!75!black, attach title to upper,
            phantom={\setcounter{myrasternum}{\thetcbrasternum}},
            title={(\alph{myrasternum})\enspace}, halign lower=center, halign
            upper = left, lower separated=false, middle=0mm, boxsep=0mm,
            leftupper=2mm, left=2mm, right=2mm, fonttitle=\bfseries,
            coltitle=black, raster column skip=1mm,raster halign=center, raster
            row skip=1mm]

            \tcbitem
                \tcblower
                    \scalebox{0.4}{
                        \begin{tikzpicture}
                            \petersentres
                            \unopBunopRp{2.5cm}{90 + (360/5)*2}
                            \unotGunotYt{5cm}{90 + (360/5)}

                        \end{tikzpicture}
                    }
            
            \tcbitem
                \tcblower
                    \scalebox{0.4}{
                        \begin{tikzpicture}
                            \petersencuatro
                            \unopBunopRp{2.5cm}{90 + (360/5)*2}
                            \unotGunotYt{5cm}{90 + (360/5)}
                            \unopGunopY{5cm}{90 + (360/5)*3}
                            
                        \end{tikzpicture}
                    }
            \tcbitem
                \tcblower
                    \scalebox{0.4}{
                        \begin{tikzpicture}
                            \petersen
                            \unopBunopRp{2.5cm}{90 + (360/5)*2}
                            \unotGunotYt{5cm}{90 + (360/5)}
                            \unopG{2.5cm}{90 + (360/5)*3}
                            \unopY{5cm}{90 + (360/5)*4}
                            
                        \end{tikzpicture}
                    }
        \end{tcbitemize}
        \captionsetup{font=scriptsize}
        \caption{En los siguientes diagramas, las fichas $c_{1,2}$ y $c_{3,2}$
        se omiten. En (a), $R$ acaba de mover $r_1$, de modo que $c_{1,1}$ y
        $c_{3,1}$ le han recapturado. Luego, se busca posicionar a $c_{2,1}$ y
        $c_{4,1}$ adyacentes a $u_1$ y $u_2$. Si alguno entre $c_{2,1}$ y
        $c_{4,1}$ estuviera en un v\'ertice negro sombreado, $C$ ganar\'ia en a
        lo m\'as un turno o un polic\'ia tendr\'ia sus dos fichas en un mismo
        v\'ertice, lo cual es imposible. De forma similar, si alguno entre
        $c_{2,1}$ y $c_{4,1}$ estuviera en $u_1$ o en $u_2$, entonces $C$ gana
        en su siguiente turno. Finalmente, para cada uno de los v\'ertices
        restantes, sombreados o no, existe un v\'ertice azul sombreado y un
        v\'ertice verde sombreado en su vecindad cerrada, determinando el
        movimiento de ambos $c_{2,1}$ y $c_{4,1}$ como se ilustra en (b) y (c).}
        \label{fig:guarding0}
    \end{figure}

    Obs\'ervese que $r_1$ no se puede mover ni a $u_1$ ni a $u_2$ pues ello
    llevar\'ia que $R$ pierda al siguiente turno. As\'i, $R$ pasa o mueve $r_2$.
    En el primer caso, ambos $c_{2,1}$ y $c_{4,1}$ avanzan a $u_1$ y $u_2$
    respectivamente para forzar a $R$ a mover $r_1$ hacia $r_2$, o a mover a
    $r_2$. Por tanto, en cualquier caso $R$ mueve $r_2$ y se puede asumir sin
    p\'erdida de generalidad que su nueva posici\'on es $w_1$. Sea
    $N(w_1)=\l{u_1, r_2(t),y}$. Como se mencion\'o anteriormente, $c_{2,2}$ y
    $c_{4,2}$ recapturan $r_2$. Se generan dos casos dependiendo de si se tiene
    una obstrucci\'on:
    
    \begin{enumerate}
        \item Si existe una obstrucci\'on, solo uno entre $c_2$ y $c_4$ est\'a
        obstru\'ido pues $c_{2,1}$ y $c_{4,1}$ est\'an en distintos v\'ertices.
        N\'otese que la obstrucci\'on se puede resolver con un reetiquetado de
        las fichas, preservando la condici\'on de que $r_1$ no puede moverse a
        ning\'un v\'ertice m\'as que para acercarse a $r_2$. Adem\'as, al igual
        que en el caso anterior, $c_1$ y $c_3$ pueden acercarse a la vecindad de
        $r_2(t+1)$, de manera que $R$ no puede mantener sus fichas en v\'ertices
        no adyacentes (v\'ease \cref{fig:guarding}).
    
        \begin{figure}[!htbp]
            \begin{tcbitemize}[raster columns=2, raster equal height, enhanced,
                colback=Black!5!white,colframe=Grey!75!black, attach title to upper,
                phantom={\setcounter{myrasternum}{\thetcbrasternum}},
                title={(\alph{myrasternum})\enspace}, halign lower=center, halign
                upper = left, lower separated=false, middle=0mm, boxsep=0mm,
                leftupper=2mm, left=2mm, right=2mm, fonttitle=\bfseries,
                coltitle=black, raster column skip=1mm,raster halign=center, raster
                row skip=1mm]
    
                \tcbitem
                    \tcblower
                        \scalebox{0.4}{
                            \begin{tikzpicture}
                                \petersencinco
                                \unopBunopRp{2.5cm}{90 + (360/5)*2}
                                \unotGunotY{5cm}{90 + (360/5)}
                                \unopGt{5cm}{90}
                                \unotB{2.5cm}{90 + (360/5)*4}
                                \unotR{2.5cm}{90 + (360/5)*3}
                                \unopY{5cm}{90 + (360/5)*4}
    
                            \end{tikzpicture}
                        }
                
                \tcbitem
                    \tcblower
                        \scalebox{0.4}{
                            \begin{tikzpicture}
                                \petersen
                                \unopBunopRp{2.5cm}{90 + (360/5)*2}
                                \unopG{5cm}{90 + (360/5)}
                                \unotGunotYt{5cm}{90}
                                \unotB{2.5cm}{90 + (360/5)*4}
                                \unotR{2.5cm}{90 + (360/5)}
                                \unopY{5cm}{90 + (360/5)*4}
                                
                            \end{tikzpicture}
                    }
                \tcbitem
                    \tcblower
                        \scalebox{0.4}{
                            \begin{tikzpicture}
                                \petersensiete
                                \unopBunopRp{2.5cm}{90 + (360/5)*2}
                                \unopG{5cm}{90 + (360/5)}
                                \unotGunotYt{5cm}{90}
                                \unotB{2.5cm}{90 + (360/5)*4}
                                \unotR{2.5cm}{90 + (360/5)*1}
                                \unopY{5cm}{90 + (360/5)*4}
                                
                            \end{tikzpicture}
                        }
                \tcbitem
                    \tcblower
                        \scalebox{0.4}{
                            \begin{tikzpicture}
                                \petersen
                                \unopBunopRp{2.5cm}{90 + (360/5)*2}
                                \unopG{5cm}{90 + (360/5)*2}
                                \unotGunotYt{5cm}{90}
                                \unotB{5cm}{90 + (360/5)*4}
                                \unotR{5cm}{90 + (360/5)*1}
                                \unopY{2.5cm}{90 + (360/5)*4}
                                
                            \end{tikzpicture}
                        }
            \end{tcbitemize}
            \captionsetup{font=scriptsize}
            \caption{En (a), $R$ acaba de mover $r_2$, de manera que $c_{2,2}$ y
            $c_{4,2}$ recapturan y $c_2$ reetiqueta sus fichas como en (b).
            Mientras tanto, $c_{1,2}$ y $c_{3,2}$ se mueven como en (b) para
            posicionarse en un v\'ertice adyacente a $y$ y un v\'ertice
            adyacente a $r_2(t)$ respectivamente. N\'otese que el reetiquetado
            causa que $c_{2,1}$ sea adyacente al v\'ertice azul sin sombreado en
            (c), de modo que este a\'un cumple su funci\'on del paso anterior.
            V\'ease que $R$ solo puede pasar si es que quiere evitar que $r_1$ y
            $r_2$ est\'en en v\'ertices adyacentes. No obstante, $C$ mueve sus
            polic\'ias como en (d), forzando el movimiento de $R$.}
            \label{fig:guarding}
        \end{figure}

        \item Si no se tiene una obstrucci\'on, ambos $c_{1,2}$ y
        $c_{3,2}$ puede posicionarse en la vecindad de $y$. Aqu\'i se generan
        tres subcasos:
        \begin{enumerate}
            \item Si alguno entre $c_{1,2}$ y $c_{3,2}$ est\'a en $w_2$ o en
            $v$ (v\'ease (b) en \cref{fig:guarding1}), sup\'ongase sin p\'erdida
            de generalidad que es $c_{1,2}$, de modo que este se mueve hacia
            $r_2(t)$ y $c_{3,2}$ se posiciona en la vecindad de $y$ (v\'ease (c)
            en \cref{fig:guarding1}). Esto obliga a $R$ a mover $r_1$ a $v$.

            \begin{figure}[!htbp]
                \begin{tcbitemize}[raster columns=2, raster equal height, enhanced,
                    colback=Black!5!white,colframe=Grey!75!black, attach title to upper,
                    phantom={\setcounter{myrasternum}{\thetcbrasternum}},
                    title={(\alph{myrasternum})\enspace}, halign lower=center, halign
                    upper = left, lower separated=false, middle=0mm, boxsep=0mm,
                    leftupper=2mm, left=2mm, right=2mm, fonttitle=\bfseries,
                    coltitle=black, raster column skip=1mm,raster halign=center, raster
                    row skip=1mm]
    
                    \tcbitem
                        \tcblower
                            \scalebox{0.4}{
                                \begin{tikzpicture}
                                    \petersen
                                    \unopBunopRp{2.5cm}{90 + (360/5)*2}
                                    \unotGunotY{5cm}{90 + (360/5)}
                                    \tri{5cm}{90}
                                    \unotBunotR{5cm}{90 + (360/5)*2}
                                    \unopG{2.5cm}{90 + (360/5)*3}
                                    \unopY{5cm}{90 + (360/5)*4}
    
                                \end{tikzpicture}
                            }
                    
                    \tcbitem
                        \tcblower
                            \scalebox{0.4}{
                                \begin{tikzpicture}
                                    \petersenseis
                                    \unopBunopRp{2.5cm}{90 + (360/5)*2}
                                    \unotGunotYt{5cm}{90}
                                    \unotBunotR{5cm}{90 + (360/5)*2}
                                    \unopG{2.5cm}{90 + (360/5)*3}
                                    \unopY{5cm}{90 + (360/5)*4}
                                    
                                \end{tikzpicture}
                            }
                    \tcbitem[raster multicolumn=2]
                        \tcblower
                            \scalebox{0.4}{
                                \begin{tikzpicture}
                                    \petersen
                                    \unopBunopRp{2.5cm}{90 + (360/5)*2}
                                    \unotGunotYt{5cm}{90}
                                    \unotR{5cm}{90 + (360/5)*3}
                                    \unotB{5cm}{90 + (360/5)}
                                    \unopG{2.5cm}{90 + (360/5)*3}
                                    \unopY{5cm}{90 + (360/5)*4}
                                    
                                \end{tikzpicture}
                            }
                \end{tcbitemize}
                \captionsetup{font=scriptsize}
                \caption{Cuando el movimiento de $r_2$ no propicie una
                obstrucci\'on de $c_2$ o $c_4$, ambos polic\'ias recapturan como
                es usual, mientras que $c_{1,2}$ se mueve hacia $r_2(t)$ y
                $c_{3,2}$ se posiciona en la vecindad de $y$. N\'otese que se
                ignorar\'an los casos en los que alguno entre $c_{1,2}$ y
                $c_{3,2}$ est\'an en los v\'ertices sombreados de negro o en $y$
                pues $C$ gana en a lo m\'as un turno o est\'a en una posici\'on
                ilegal.}
                \label{fig:guarding1}
            \end{figure}

            \item Si ninguno est\'a en $w_2$ o $v$, pero al menos uno est\'a en
            $u_2$ o $z$ y el otro no est\'a en $z$, sup\'ongase sin p\'erdida de
            generalidad que $c_{1,2}$ est\'a en $u_2$. As\'i, $c_{1,2}$ se mueve
            a $w_2$ y, como $c_{3,2}$ no est\'a en $z$, este se mueve a $y$,
            obligando nuevamente a que $R$ mueva $r_1$ a $v$ (v\'ease
            \cref{fig:guarding2}).

            \begin{figure}[!htbp]
                \begin{tcbitemize}[raster columns=2, raster equal height, enhanced,
                    colback=Black!5!white,colframe=Grey!75!black, attach title to upper,
                    phantom={\setcounter{myrasternum}{\thetcbrasternum}},
                    title={(\alph{myrasternum})\enspace}, halign lower=center, halign
                    upper = left, lower separated=false, middle=0mm, boxsep=0mm,
                    leftupper=2mm, left=2mm, right=2mm, fonttitle=\bfseries,
                    coltitle=black, raster column skip=1mm,raster halign=center, raster
                    row skip=1mm]
    
                    \tcbitem
                        \tcblower
                            \scalebox{0.4}{
                                \begin{tikzpicture}
                                    \petersen
                                    \unopBunopRp{2.5cm}{90 + (360/5)*2}
                                    \unotGunotY{5cm}{90 + (360/5)}
                                    \tri{5cm}{90}
                                    \unotB{2.5cm}{90 + (360/5)*4}
                                    \unotR{5cm}{90 + (360/5)*3}
                                    \unopG{2.5cm}{90 + (360/5)*3}
                                    \unopY{5cm}{90 + (360/5)*4}
    
                                \end{tikzpicture}
                            }
                    
                    \tcbitem
                        \tcblower
                            \scalebox{0.4}{
                                \begin{tikzpicture}
                                    \petersenseis
                                    \unopBunopRp{2.5cm}{90 + (360/5)*2}
                                    \unotGunotYt{5cm}{90}
                                    \unotB{2.5cm}{90 + (360/5)*4}
                                    \unotR{5cm}{90 + (360/5)*3}
                                    \unopG{2.5cm}{90 + (360/5)*3}
                                    \unopY{5cm}{90 + (360/5)*4}
                                    
                                \end{tikzpicture}
                            }
                    \tcbitem[raster multicolumn=2]
                        \tcblower
                            \scalebox{0.4}{
                                \begin{tikzpicture}
                                    \petersen
                                    \unopBunopRp{2.5cm}{90 + (360/5)*2}
                                    \unotGunotYt{5cm}{90}
                                    \unotB{2.5cm}{90 + (360/5)}
                                    \unopG{2.5cm}{90 + (360/5)*3}
                                    \unotRunopY{5cm}{90 + (360/5)*4}
                                    
                                \end{tikzpicture}
                            }
                \end{tcbitemize}
                \captionsetup{font=scriptsize}
                \caption{Al igual que en el subcaso anterior, $c_2$ y $c_4$
                recapturan como en (b), mientras que $c_{1,2}$ y $c_{3,2}$ se
                mueven como en (c) para obligar el movimiento de $r_1$ a $v$.}
                \label{fig:guarding2}
            \end{figure}

            \item Si los anteriores casos no suceden, $c_{1,2}$ y $c_{3,2}$
            deben estar ambos en $z$ o en $x$ (v\'ease (b) en
            \cref{fig:guarding3}). En este caso, $C$ mueve las fichas de manera
            que $c_{1,2}$ termina en $z$ y $c_{3,2}$ est\'a en $x$.

            \begin{figure}[!htbp]
                \begin{tcbitemize}[raster columns=2, raster equal height, enhanced,
                    colback=Black!5!white,colframe=Grey!75!black, attach title to upper,
                    phantom={\setcounter{myrasternum}{\thetcbrasternum}},
                    title={(\alph{myrasternum})\enspace}, halign lower=center, halign
                    upper = left, lower separated=false, middle=0mm, boxsep=0mm,
                    leftupper=2mm, left=2mm, right=2mm, fonttitle=\bfseries,
                    coltitle=black, raster column skip=1mm,raster halign=center, raster
                    row skip=1mm]
    
                    \tcbitem
                        \tcblower
                            \scalebox{0.4}{
                                \begin{tikzpicture}
                                    \petersen
                                    \unopBunopRp{2.5cm}{90 + (360/5)*2}
                                    \unotGunotY{5cm}{90 + (360/5)}
                                    \tri{5cm}{90}
                                    \unotBunotR{5cm}{90 + (360/5)*3}
                                    \unopG{2.5cm}{90 + (360/5)*3}
                                    \unopY{5cm}{90 + (360/5)*4}
    
                                \end{tikzpicture}
                            }
                    
                    \tcbitem
                        \tcblower
                            \scalebox{0.4}{
                                \begin{tikzpicture}
                                    \petersenocho
                                    \unopBunopRp{2.5cm}{90 + (360/5)*2}
                                    \unotGunotYt{5cm}{90}
                                    \unotBunotR{5cm}{90 + (360/5)*3}
                                    \unopG{2.5cm}{90 + (360/5)*3}
                                    \unopY{5cm}{90 + (360/5)*4}
                                    
                                \end{tikzpicture}
                            }
                    \tcbitem[raster multicolumn=2]
                        \tcblower
                            \scalebox{0.4}{
                                \begin{tikzpicture}
                                    \petersen
                                    \unopBunopRp{2.5cm}{90 + (360/5)*2}
                                    \unotGunotYt{5cm}{90}
                                    \unotBunopG{2.5cm}{90 + (360/5)*3}
                                    \unopY{5cm}{90 + (360/5)*4}
                                    \unotR{5cm}{90 + (360/5)*3}
                                    
                                \end{tikzpicture}
                            }
                \end{tcbitemize}
                \captionsetup{font=scriptsize}
                \caption{Una vez m\'as, $c_2$ y $c_4$ recapturan como en (b),
                mientras que $c_{1,2}$ se posiciona en $z$ y $c_{3,2}$ en $x$.}
                \label{fig:guarding3}
            \end{figure}

            Si $R$ pasa, la posici\'on actual es una an\'aloga a la mostrada en
            (b) de \cref{fig:guarding2}, de modo que $C$ mueve como en el caso
            b) para obligar el movimiento de $r_1$ a $v$. De lo contrario, $R$
            mueve $r_2$ a $r_2(t)$, de manera que $c_{2,2}$ y $c_{4,2}$ le
            recapturan y $c_{1,2}$ se mueve a $w_2$. En esta posici\'on $R$ no
            puede mover $r_1$ a $v$ pues pondr\'ia sus fichas en v\'ertices
            adyacentes, de modo que $r_2$ regresa a $w_1$. As\'i $C$ est\'a en
            una posici\'on an\'aloga a la mostrada en (b) de
            \cref{fig:guarding1}, por lo que puede moverse como en el caso a)
            para obligar el movimiento de $r_1$ a $v$.

        \end{enumerate}
        
        Por tanto, en cualquier caso $R$ est\'a forzado a mover $r_1$ hacia $v$.
        Al igual que antes, $c_{1,1}$ y $c_{3,1}$ recapturan $r_1$, mientras
        que, sin p\'erdida de generalidad, $c_{2,1}$ y $c_{4,1}$ se posicionan
        en $x$ y $u_2$ respectivamente. M\'as a\'un, $R$ est\'a en una
        posici\'on an\'aloga alguna de las mostradas en (c) para
        \cref{fig:guarding1,fig:guarding2,fig:guarding3}. Si se viene de la
        posici\'on (c) de \cref{fig:guarding3}, el \'unico movimiento que puede
        hacer $R$ es mover $r_2$ o $r_1$ hacia $r_2(t)$, poni\'endo las fichas
        de $r$ en v\'ertices adyacentes (v\'ease \cref{fig:guarding4}).
        
        \begin{figure}[!htbp]
            \begin{tcbitemize}[raster columns=2, raster equal height, enhanced,
                colback=Black!5!white,colframe=Grey!75!black, attach title to upper,
                phantom={\setcounter{myrasternum}{\thetcbrasternum}},
                title={(\alph{myrasternum})\enspace}, halign lower=center, halign
                upper = left, lower separated=false, middle=0mm, boxsep=0mm,
                leftupper=2mm, left=2mm, right=2mm, fonttitle=\bfseries,
                coltitle=black, raster column skip=1mm,raster halign=center, raster
                row skip=1mm]
    
                \tcbitem
                    \tcblower
                        \scalebox{0.4}{
                            \begin{tikzpicture}
                                \petersen
                                \unopBunopR{2.5cm}{90 + (360/5)*2}
                                \punto{5cm}{90 + (360/5)*2}
                                \unotGunotYt{5cm}{90}
                                \unotBunopG{2.5cm}{90 + (360/5)*3}
                                \unopY{5cm}{90 + (360/5)*4}
                                \unotR{5cm}{90 + (360/5)*3}
                                
                            \end{tikzpicture}
                        }
                
                \tcbitem
                    \tcblower
                        \scalebox{0.4}{
                            \begin{tikzpicture}
                                \petersen
                                \unopBunopRp{5cm}{90 + (360/5)*2}
                                \unotGunotYt{5cm}{90}
                                \unotB{2.5cm}{90 + (360/5)*3}
                                \unopGunotR{5cm}{90 + (360/5)*3}
                                \unopY{2.5cm}{90 + (360/5)*4}
                                
                            \end{tikzpicture}
                        }
            \end{tcbitemize}
            \captionsetup{font=scriptsize}
            \caption{En (a), $R$ acaba de mover $r_1$, de manera que $C$ mueve
            como en (b) para posicionarse $c_{2,1}$ en $x$ y $c_{4,1}$ en $u_2$,
            a la vez que $c_{1,1}$ y $c_{3,1}$ recapturan. En este caso, $R$ no
            tiene ning\'un movimiento mas que mover $r_1$ hacia $r_2$ o
            viceversa.}
            \label{fig:guarding4}
        \end{figure}

        Si se viene de la posici\'on (c) de \cref{fig:guarding2}, $R$ puede
        mover $r_2$ a $u_1$. No obstante, si ambos $c_{2,2}$ y $c_{4,2}$
        recapturan, mientras que $c_{1,2}$ se mueve a $z$, $R$ se ve obligado a
        poner $r_1$ y $r_2$ en v\'ertices adyacentes (v\'ease
        \cref{fig:guarding5}).

        \begin{figure}[!htbp]
            \begin{tcbitemize}[raster columns=2, raster equal height, enhanced,
                colback=Black!5!white,colframe=Grey!75!black, attach title to upper,
                phantom={\setcounter{myrasternum}{\thetcbrasternum}},
                title={(\alph{myrasternum})\enspace}, halign lower=center, halign
                upper = left, lower separated=false, middle=0mm, boxsep=0mm,
                leftupper=2mm, left=2mm, right=2mm, fonttitle=\bfseries,
                coltitle=black, raster column skip=1mm,raster halign=center, raster
                row skip=1mm]
    
                \tcbitem
                    \tcblower
                        \scalebox{0.4}{
                            \begin{tikzpicture}
                                \petersen
                                \unopBunopRp{5cm}{90 + (360/5)*2}
                                \unotGunotY{5cm}{90}
                                \tri{2.5cm}{90}
                                \unotB{2.5cm}{90 + (360/5)}
                                \unopG{5cm}{90 + (360/5)*3}
                                \unotR{5cm}{90 + (360/5)*4}
                                \unopY{2.5cm}{90 + (360/5)*4}
                                
                            \end{tikzpicture}
                        }
                \tcbitem
                    \tcblower
                        \scalebox{0.4}{
                            \begin{tikzpicture}
                                \petersen
                                \unopBunopRp{5cm}{90 + (360/5)*2}
                                \unotGunotYt{2.5cm}{90}
                                \unotB{2.5cm}{90 + (360/5)*3}
                                \unopG{5cm}{90 + (360/5)*3}
                                \unotR{5cm}{90 + (360/5)*4}
                                \unopY{2.5cm}{90 + (360/5)*4}
                                
                            \end{tikzpicture}
                        }
            \end{tcbitemize}
            \captionsetup{font=scriptsize}
            \caption{Si $R$ mueve $r_2$ a $u_2$ como en (a), $C$ mueve a sus
            polic\'ias como en (b) para forzar el movimiento de $R$ que ponga a
            $r_1$ adyacente a $r_2$.}
            \label{fig:guarding5}
        \end{figure}

        Finalmente, si se viene de la posici\'on (c) de \cref{fig:guarding1},
        $R$ puede mover $r_2$ a $u_1$. Esta situaci\'on tiene de hecho dos
        variantes, dependiendo de si $c_{3,2}$ est\'a en $u_2$ o en $x$. No
        obstante, en cualquiera de los casos se tiene que $c_{3,2}$ est\'a en un
        v\'ertice adyacente a un v\'ertice de la vecindad cerrada de $z$. As\'i,
        en su turno, $C$ recaptura con $c_{2,2}$ y $c_{4,2}$ y $c_{3,2}$ se
        acerca a $z$. Por tanto, $R$ no puede mover $r_2$ a $z$ pues este pierde
        en a lo m\'as un turno, por lo que se ve obligado a juntar las fichas
        $r_1$ y $r_2$ en v\'ertices adyacentes (v\'ease \cref{fig:guarding6}).

        \begin{figure}[!htbp]
            \begin{tcbitemize}[raster columns=2, raster equal height, enhanced,
                colback=Black!5!white,colframe=Grey!75!black, attach title to upper,
                phantom={\setcounter{myrasternum}{\thetcbrasternum}},
                title={(\alph{myrasternum})\enspace}, halign lower=center, halign
                upper = left, lower separated=false, middle=0mm, boxsep=0mm,
                leftupper=2mm, left=2mm, right=2mm, fonttitle=\bfseries,
                coltitle=black, raster column skip=1mm,raster halign=center, raster
                row skip=1mm]
    
                \tcbitem
                    \tcblower
                        \scalebox{0.4}{
                            \begin{tikzpicture}
                                \petersen
                                \unopBunopRp{5cm}{90 + (360/5)*2}
                                \unotGunotY{5cm}{90}
                                \tri{2.5cm}{90}
                                \unotRunopY{2.5cm}{90 + (360/5)*4}
                                \unotB{5cm}{90 + (360/5)}
                                \unopG{5cm}{90 + (360/5)*3}
                                
                            \end{tikzpicture}
                        }
                \tcbitem
                    \tcblower
                        \scalebox{0.4}{
                            \begin{tikzpicture}
                                \petersen
                                \unopBunopRp{5cm}{90 + (360/5)*2}
                                \unotGunotYt{2.5cm}{90}
                                \unopY{2.5cm}{90 + (360/5)*4}
                                \unotR{2.5cm}{90+ (360/5)*1}
                                \unotB{5cm}{90 + (360/5)}
                                \unopG{5cm}{90 + (360/5)*3}
                                
                            \end{tikzpicture}
                        }
            \end{tcbitemize}
            \captionsetup{font=scriptsize}
            \caption{En (a), $R$ acaba de mover $r_2$ a $u_1$, por lo que $C$
            mueve como en (b) para forzar el movimiento de $r_1$ hacia $r_2$
            o viceversa.}
            \label{fig:guarding6}
        \end{figure}
    \end{enumerate}

    Esto concluye la prueba de que $r_1$ y $r_2$ eventualmente est\'an en
    posiciones adyacentes. M\'as a\'un, de las escenas (b) de
    \cref{fig:guarding4,fig:guarding5,fig:guarding6}, de las escenas (c) de
    \cref{fig:guarding0,,fig:guarding1,fig:guarding2,fig:guarding3}, y de las
    escenas (c) y (d) de \cref{fig:guarding}, que antes de que $r_1$ y $r_2$
    est\'en en v\'ertices adyacentes, $c_{1,1}$ y $c_{3,1}$ capturaban a $r_1$,
    mientras que $c_{2,2}$ y $c_{4,2}$ capturaban a $r_2$. As\'i, sup\'ongase
    sin p\'erdida de generalidad que $r_1$ fue la \'ultima ficha movida, de modo
    que la situaci\'on es como se muestra en \cref{fig:guarding7}. Los
    polic\'ias proceden a posicionarse sobre las vecindades de los v\'ertices a
    los que puede mover $R$, dejando a $r_1$ sin movimientos.
    
    \begin{figure}[!htbp]
        \begin{tcbitemize}[raster columns=2, raster equal height, enhanced,
            colback=Black!5!white,colframe=Grey!75!black, attach title to upper,
            phantom={\setcounter{myrasternum}{\thetcbrasternum}},
            title={(\alph{myrasternum})\enspace}, halign lower=center, halign
            upper = left, lower separated=false, middle=0mm, boxsep=0mm,
            leftupper=2mm, left=2mm, right=2mm, fonttitle=\bfseries,
            coltitle=black, raster column skip=1mm,raster halign=center, raster
            row skip=1mm]
            \tcbitem
                \tcblower
                    \scalebox{0.4}{
                        \begin{tikzpicture}
                            \petersennueve
                            \punto{5cm}{90 + (360/5)}
                            \unopBunopR{5cm}{90 + (360/5)*2}
                            \unotGunotYt{5cm}{90}
                            
                        \end{tikzpicture}
                    }
            
            \tcbitem
                \tcblower
                    \scalebox{0.4}{
                        \begin{tikzpicture}
                            \petersendiez
                            \punto{5cm}{90 + (360/5)}
                            \unopBunopR{5cm}{90 + (360/5)*2}
                            \unotGunotYt{5cm}{90}
                            \unopG{2.5cm}{90 + (360/5)*2}
                            \unopY{2.5cm}{90 + (360/5)*3}
                        \end{tikzpicture}
                    }
            
            \tcbitem[raster multicolumn=2]
                \tcblower
                    \scalebox{0.4}{
                        \begin{tikzpicture}
                            \petersen
                            \punto{5cm}{90 + (360/5)}
                            \unopBunopR{5cm}{90 + (360/5)*2}
                            \unotGunotYt{5cm}{90}
                            \unotB{2.5cm}{90 + (360/5)*4}
                            \unopG{2.5cm}{90 + (360/5)*2}
                            \unotRunopY{2.5cm}{90 + (360/5)*3}
                        \end{tikzpicture}
                    }
        \end{tcbitemize}
        \captionsetup{font=scriptsize}
        \caption{En (a), $C$ busca poner, sin p\'erdida de generalidad, a
        $c_{2,1}$ y $c_{4,1}$ en un v\'ertice azul sombreado y en un v\'ertice
        sombreado verde. Si alguno de ellos estuviera en los v\'ertices negros
        sombreados, en el v\'ertice azul no sombreado, o en el v\'ertice verde
        no sombreado, se sigue que $C$ captura a $R$ en a lo m\'as un turno o
        est\'a en una configuraci\'on inicial. Por otra parte el resto de
        v\'ertices es tal que tiene al menos un v\'ertice verde sombreado y uno
        azul sombreado en su vecindad cerrada. Por tanto, siempre se pueden
        colocar $c_{2,1}$ y $c_{4,1}$ de manera que estos impiden que $r_1$ se
        mueva los v\'ertices verde o azul sin sombrear. An\'alogamente se mueven
        $c_{1,2}$ y $c_{3,2}$ en un v\'ertice azul y uno verde en (b). As\'i,
        $C$ termina su turno como se muestra en (c).}
        \label{fig:guarding7}
    \end{figure}

    Por \'ultimo, si $R$ pasa, $C$ mueve como en \cref{fig:captura1} para
    capturar a $R$ en su siguiente turno.

    \begin{figure}[!htbp]
        \centering
        \scalebox{0.4}{
            \begin{tikzpicture}
                \petersen
                \unopBunopRp{5cm}{90 + (360/5)}
                \unotGunotYt{5cm}{90}
                \unotB{2.5cm}{90 + (360/5)*4}
                \unopG{5cm}{90 + (360/5)*2}
                \unotR{2.5cm}{90 + (360/5)*3}
                \unopY{2.5cm}{90 + (360/5)*1}
            \end{tikzpicture}
        }
    \captionsetup{font=scriptsize}
    \caption{Posterior a que $R$ pase, $C$ mueve para evitar que este mueva y
    capturarle en su siguiente turno.}
    \label{fig:captura1}
    \end{figure}

    Si en cambio $R$ mueve $r_2$, sup\'ongase sin p\'erdida de generalidad que
    $r_2$ se mueve a $u_1$, el otro caso es an\'alogo. Luego $C$ mueve sus
    fichas como en \cref{fig:captura2}, capturando a $R$ en su siguiente turno.

    \begin{figure}[!htbp]
        \centering
        \scalebox{0.4}{
            \begin{tikzpicture}
                \petersen
                \unopBunopRp{5cm}{90 + (360/5)}
                \unotGunotY{5cm}{90}
                \tri{2.5cm}{90}
                \unotB{2.5cm}{90 + (360/5)*4}
                \unopG{5cm}{90 + (360/5)*2}
                \unotR{2.5cm}{90 + (360/5)*3}
                \unopY{2.5cm}{90 + (360/5)*1}
            \end{tikzpicture}
        }
    \captionsetup{font=scriptsize}
    \caption{Posterior a que $R$ pase, $C$ mueve para evitar que este mueva y
    capturarle en su siguiente turno.}
    \label{fig:captura2}
    \end{figure}

    Se concluye que $c(\t{2}{P})\leq 4$.

\end{proof}